%-----------------------------------------------------------------------------------------------------------------------------------------------%
%	The MIT License (MIT)
%
%	Copyright (c) 2021 Jitin Nair
%
%	Permission is hereby granted, free of charge, to any person obtaining a copy
%	of this software and associated documentation files (the "Software"), to deal
%	in the Software without restriction, including without limitation the rights
%	to use, copy, modify, merge, publish, distribute, sublicense, and/or sell
%	copies of the Software, and to permit persons to whom the Software is
%	furnished to do so, subject to the following conditions:
%	
%	THE SOFTWARE IS PROVIDED "AS IS", WITHOUT WARRANTY OF ANY KIND, EXPRESS OR
%	IMPLIED, INCLUDING BUT NOT LIMITED TO THE WARRANTIES OF MERCHANTABILITY,
%	FITNESS FOR A PARTICULAR PURPOSE AND NONINFRINGEMENT. IN NO EVENT SHALL THE
%	AUTHORS OR COPYRIGHT HOLDERS BE LIABLE FOR ANY CLAIM, DAMAGES OR OTHER
%	LIABILITY, WHETHER IN AN ACTION OF CONTRACT, TORT OR OTHERWISE, ARISING FROM,
%	OUT OF OR IN CONNECTION WITH THE SOFTWARE OR THE USE OR OTHER DEALINGS IN
%	THE SOFTWARE.
%	
%
%-----------------------------------------------------------------------------------------------------------------------------------------------%

%----------------------------------------------------------------------------------------
%	DOCUMENT DEFINITION
%----------------------------------------------------------------------------------------

% article class because we want to fully customize the page and not use a cv template
\documentclass[a4paper,12pt]{article}

%----------------------------------------------------------------------------------------
%	FONT
%----------------------------------------------------------------------------------------

% % fontspec allows you to use TTF/OTF fonts directly
% \usepackage{fontspec}
% \defaultfontfeatures{Ligatures=TeX}

% % modified for ShareLaTeX use
% \setmainfont[
% SmallCapsFont = Fontin-SmallCaps.otf,
% BoldFont = Fontin-Bold.otf,
% ItalicFont = Fontin-Italic.otf
% ]
% {Fontin.otf}

%----------------------------------------------------------------------------------------
%	PACKAGES
%----------------------------------------------------------------------------------------
\usepackage{url}
\usepackage{parskip} 	

%other packages for formatting
\RequirePackage{color}
\RequirePackage{graphicx}
\usepackage[usenames,dvipsnames]{xcolor}
\usepackage[scale=0.92]{geometry}

%tabularx environment
\usepackage{tabularx}

%for lists within experience section
\usepackage{enumitem}

% centered version of 'X' col. type
\newcolumntype{C}{>{\centering\arraybackslash}X} 

%to prevent spillover of tabular into next pages
\usepackage{supertabular}
\usepackage{tabularx}
\newlength{\fullcollw}
\setlength{\fullcollw}{0.47\textwidth}

%custom \section
\usepackage{titlesec}				
\usepackage{multicol}
\usepackage{multirow}

%CV Sections inspired by: 
%http://stefano.italians.nl/archives/26
\titleformat{\section}{\Large\scshape\raggedright}{}{0em}{}[\titlerule]
\titlespacing{\section}{0pt}{9pt}{7pt}

%for publications
\usepackage[style=authoryear,sorting=ynt, maxbibnames=2]{biblatex}

%Setup hyperref package, and colours for links
\usepackage[unicode, draft=false]{hyperref}
\definecolor{linkcolour}{rgb}{0,0.2,0.6}
\hypersetup{colorlinks,breaklinks,urlcolor=linkcolour,linkcolor=linkcolour}
\addbibresource{citations.bib}
\setlength\bibitemsep{1em}

%for social icons
\usepackage{fontawesome5}

%debug page outer frames
%\usepackage{showframe}


% job listing environments
\newenvironment{jobshort}[2]
    {
    \begin{tabularx}{\linewidth}{@{}X r@{}}
    \textbf{#1} & #2 \\[3.75pt]
    \end{tabularx}
    }
    {
    }

\newenvironment{joblong}[2]
    {
    \begin{tabularx}{\linewidth}{@{}X r@{}}
    \textbf{#1} & #2 \\[3.75pt]
    \end{tabularx}
    \begin{minipage}[t]{\linewidth}
    \begin{itemize}[nosep,after=\strut, leftmargin=1em, itemsep=3pt,label=--]
    }
    {
    \end{itemize}
    \end{minipage}    
    }



%----------------------------------------------------------------------------------------
%	BEGIN DOCUMENT
%----------------------------------------------------------------------------------------
\begin{document}

% non-numbered pages
\pagestyle{empty} 

%----------------------------------------------------------------------------------------
%	TITLE
%----------------------------------------------------------------------------------------

% \begin{tabularx}{\linewidth}{ @{}X X@{} }
% \huge{Your Name}\vspace{2pt} & \hfill \emoji{incoming-envelope} email@email.com \\
% \raisebox{-0.05\height}\faGithub\ username \ | \
% \raisebox{-0.00\height}\faLinkedin\ username \ | \ \raisebox{-0.05\height}\faGlobe \ mysite.com  & \hfill \emoji{calling} number
% \end{tabularx}

\begin{tabularx}{\linewidth}{@{}X@{}}
\Huge{Han Li} \\[3pt]
\href{mailto:231220161@smail.nju.edu.cn}{231220161@smail.nju.edu.cn}
\ $|$ \
\href{tel:+8617391928803}{+86-173-9192-8803}
\ $|$ \
\href{https://github.com/schwerli}{GitHub: schwerli}
\ $|$ \
\href{https://scholar.google.com/citations?user=qbDL1kwAAAAJ\&hl=en}{Google Scholar}
\ $|$ \
\href{https://han-li-research.github.io}{Homepage} \\
\end{tabularx}
\vspace{0.2cm}

%----------------------------------------------------------------------------------------
% EXPERIENCE SECTIONS
%----------------------------------------------------------------------------------------

%----------------------------------------------------------------------------------------
%	EDUCATION
%----------------------------------------------------------------------------------------
\section{Education}
\begin{tabularx}{\linewidth}{@{}l X@{}} 	
2023 - present & Bachelor's Degree in Computer Science and Technology at \textbf{Nanjing University} \\
\end{tabularx}

%----------------------------------------------------------------------------------------
%	SUMMARY
%----------------------------------------------------------------------------------------
\section{Summary}
Undergraduate researcher in computer science at Nanjing University, currently interning at Microsoft Research Asia (DKI), Beijing. My recent work focuses on AI for software engineering, especially code agents, long-horizon evaluation, context retrieval, and trajectory tracing/debugging. I am primarily interested in building benchmarks and analysis tools that make agent behavior more measurable, interpretable, and reliable.

%----------------------------------------------------------------------------------------
%	RESEARCH INTERESTS
%----------------------------------------------------------------------------------------
\section{Research Interests}
AI for Software Engineering, Code Agents, Long-Horizon Evaluation, LLMs for Code

%----------------------------------------------------------------------------------------
%	SKILLS
%----------------------------------------------------------------------------------------
\section{Skills}
\textbf{Programming Languages:} C/C++, Python, Java \\
\textbf{Research Areas:} Code Agents, Benchmarking, Program Analysis, Machine Learning

%----------------------------------------------------------------------------------------
%	WORK EXPERIENCE
%----------------------------------------------------------------------------------------
\section{Work Experience}
\vspace{-0.1cm}

\begin{jobshort}{Research Intern, Microsoft Research Asia (DKI), Beijing}{2026 - present}
\end{jobshort}
\vspace{-0.6cm}

\begin{jobshort}{Research Assistant, University College London}{2025.8 - present}
\end{jobshort}
\vspace{-0.6cm}

\begin{jobshort}{Research Intern, Kuaishou Core Algorithm Group}{2025.9 - 2025.12}
\end{jobshort}
\vspace{-0.6cm}

\section{Service}
\vspace{-0.1cm}

\begin{jobshort}{External Reviewer, FSE (Foundations of Software Engineering)}{2025}
\end{jobshort}

%Research Experience
\section{Research Experience}

\begin{joblong}{Long-Horizon Benchmark for Coding Agents}{2026 - present}
\item Building a benchmark for evaluating long-horizon code-agent behavior, with emphasis on multi-step planning, tool use, and trajectory-level failure analysis
\item Current project at Microsoft Research Asia (DKI), Beijing
\end{joblong}

\begin{joblong}{ContextBench: A Benchmark for Context Retrieval in Coding Agents}{2025.9 - present}
\item Developed a process-oriented benchmark for evaluating context retrieval in coding agents, with gold-context annotations and intermediate retrieval metrics
\item Constructed 1,136 issue-resolution tasks spanning 66 repositories and 8 programming languages
\item Paper: \href{https://arxiv.org/abs/2602.05892}{arXiv:2602.05892}; under review at ICML
\item \textit{Completed during research at University College London}
\end{joblong}

\begin{joblong}{CodeTracer: Towards Traceable Code Agent Trajectories}{2025.9 - 2025.12}
\item Built a scalable framework that converts raw code-agent runs into hierarchical traces and localizes the earliest failure-critical stage with compact evidence retrieval
\item Constructed CodeTraceBench from executed trajectories of four widely used code-agent frameworks over repository-level bug fixing and terminal interaction workloads
\item Joint work completed at Kuaishou in collaboration with Link Lab; ACL submission under review
\end{joblong}

\begin{joblong}{Prometheus: Multi-Agent Software Engineering System}{2025.8 - present}
\item Contributed to a multi-agent framework for repository-level software engineering tasks with strong performance on SWE-bench Verified
\item Focused on agent design and engineering for reliable code-task execution
\item Paper: \href{https://arxiv.org/abs/2507.19942}{arXiv:2507.19942}; under review at ISSTA
\item \textit{Supervised by Prof. Ye He at University College London}
\end{joblong}

\begin{joblong}{Automated Software Bug Repair with the SWE-Forge Framework}{2025.7 - present}
\item Developed SWE-Forge, a bug-fixing framework leveraging DNL model for automated root cause analysis
\item Achieved 19.8\% performance uplift on SWE-bench lite benchmark
\item \textit{Completed at Nanjing University LIP Lab}
\end{joblong}

%----------------------------------------------------------------------------------------
%	PUBLICATIONS
%----------------------------------------------------------------------------------------
\section{Selected Publications and Manuscripts}
\begin{itemize}[nosep,after=\strut, leftmargin=1em, itemsep=3pt,label=--]
\item Han Li, et al. ``ContextBench: A Benchmark for Context Retrieval in Coding Agents.'' \href{https://arxiv.org/abs/2602.05892}{arXiv:2602.05892}, under review at ICML, 2026.
\item ``CodeTracer: Towards Traceable Code Agent Trajectories.'' ACL submission under review, joint work with Kuaishou and Link Lab, 2026.
\item Yue Pan, et al. ``Prometheus: Towards Long-Horizon Codebase Navigation for Repository-Level Problem Solving.'' \href{https://arxiv.org/abs/2507.19942}{arXiv:2507.19942}, under review at ISSTA, 2026.
\item Jingxuan Xu, et al. ``SWE-Compass: Towards Unified Evaluation of Agentic Coding Abilities for Large Language Models.'' \href{https://arxiv.org/abs/2511.05459}{arXiv:2511.05459}, under review at ICML, 2025.
\item Zizheng Zhan, et al. ``KAT-Coder: A Large-Scale Agentic Code Model Trained Through Four-Stage Curriculum.'' technical report, 2025.
\end{itemize}

\vfill
\center{\footnotesize Last updated: \today}

\end{document}
